\documentclass[11pt]{article}
\usepackage[margin=1in]{geometry}
\usepackage{amsmath,amssymb,amsthm,mathtools}
\usepackage{graphicx}
\usepackage{microtype}
\usepackage[hidelinks]{hyperref}

\newtheorem{theorem}{Theorem}
\newtheorem{lemma}[theorem]{Lemma}
\newtheorem{proposition}[theorem]{Proposition}
\theoremstyle{remark}
\newtheorem{remark}[theorem]{Remark}

\newcommand{\M}{\mathrm M}
\newcommand{\id}{\mathrm{id}}
\newcommand{\Fix}{\operatorname{Fix}}
\newcommand{\K}{\mathcal K}

\title{Rieffel-Deformed Completely Positive Maps:\\
Full Answers to Questions 4.7 and 4.8}
\author{A candidate full solution to two questions in arXiv:2311.04630}
\date{23 July 2026}

\begin{document}
\maketitle

\begin{center}
\fbox{\parbox{0.91\textwidth}{\textbf{Status.}
Candidate full solution, likely valid, subject to human review.
Both questions have affirmative answers for every locally compact abelian
deformation group.}}
\end{center}

\section{Source questions}

Let $\Gamma\curvearrowright_\rho A$ be an action of a locally compact
abelian group on a $C^*$-algebra and let $\Psi$ be a continuous $2$-cocycle
on $\widehat\Gamma$. Skalski and Viselter construct, from every equivariant
nondegenerate completely positive map $S:A\to A$, a completely positive map
$S^\Psi:A^\Psi\to A^\Psi$ of norm one. They ask whether $S^\Psi$ is always
nondegenerate (Question~4.7), and whether an equivariant point-norm continuous
semigroup $(S_t)_{t\geq0}$ remains point-norm continuous after deformation
(Question~4.8)~\cite{SkalskiViselter}.

Both questions appear on PDF page~26:
\begin{figure}[ht]
  \centering
  \includegraphics[width=\textwidth]{figures/open_questions_crop.png}
  \caption{Questions 4.7--4.9 in \cite[p.~26]{SkalskiViselter}.}
\end{figure}

\noindent\textbf{Transcription of the two live questions.}
Question~4.7 asks whether $S^\Psi:A^\Psi\to A^\Psi$ is always nondegenerate.
Question~4.8 asks whether $(S_t^\Psi)_{t\geq0}$ is point-norm continuous when
$(S_t)_{t\geq0}$ is a point-norm continuous semigroup of equivariant
nondegenerate completely positive maps with $S_0=\id_A$.
Question~4.9 is not a live target here: immediately after stating it, the
source reports an affirmative solution by Shalit and Skeide.

\section{Result}

\begin{theorem}\label{thm:main}
In the full generality of Questions~4.7 and 4.8:
\begin{enumerate}
  \item $S^\Psi:A^\Psi\to A^\Psi$ is nondegenerate for every equivariant
  nondegenerate completely positive map $S:A\to A$;
  \item if $(S_t)_{t\geq0}$ is a point-norm continuous semigroup as in
  Question~4.8, then $(S_t^\Psi)_{t\geq0}$ is point-norm continuous.
\end{enumerate}
Thus both questions have affirmative answers for every locally compact
abelian deformation group $\Gamma$.
\end{theorem}

\section{Proof intuition}

Put $B=A\rtimes_\rho\Gamma$ and write
$\alpha=\widehat\rho^\Psi$ for the deformed dual action. The source proves
that the crossed-product extension $\widetilde S:B\to B$ is nondegenerate,
commutes with $\alpha$, and restricts to $S^\Psi$ on the Landstad algebra.

The strict-topology obstacle disappears after crossing one more time.
Landstad duality also writes
\[
 B=A^\Psi\rtimes_{\rho^\Psi}\Gamma,
 \qquad
 \widetilde S=S^\Psi\rtimes\Gamma .
\]
Crossing $\widetilde S$ by the dual action preserves nondegeneracy. Takai
duality identifies the resulting double-crossed map with
$S^\Psi\otimes\id_{\K(L^2\Gamma)}$. Nondegeneracy of a stabilization reflects
back to the original completely positive map. This proves Question~4.7.

For Question~4.8, the extensions $\widetilde S_t$ form a point-norm
continuous semigroup on $B$: this follows first on $C_c(\Gamma,A)$ from the
$L^1$ estimate and then by density. The Landstad algebra is the generalized
fixed-point algebra of the weakly proper
$\widehat\Gamma\rtimes\widehat\Gamma$-algebra
$(B,\alpha,\lambda)$. Its dense compactly supported core consists of
averages
\[
 E\bigl(\lambda(f)b\lambda(g)\bigr),
 \qquad f,g\in C_c(\widehat\Gamma),\ b\in B.
\]
The averaging map is norm bounded when $f$ and $g$ are fixed. Since every
$\widetilde S_t$ fixes $\lambda(C_0(\widehat\Gamma))$ and commutes with
$\alpha$, continuity on $B$ passes to this dense core, and contractivity
passes it to all of $A^\Psi$.

\section{Crossed products reflect nondegeneracy}

We use the standard crossed-product extension of an equivariant completely
positive map: if $R:C\to C$ commutes with an action
$G\curvearrowright_\beta C$, then
\[
 (R\rtimes G)(h)(s)=R(h(s)),
 \qquad h\in C_c(G,C).
\]
The source recalls that this extension exists without assuming
nondegeneracy, and that it is nondegenerate whenever $R$ is
nondegenerate~\cite[Proposition~4.5 and Remark~4.6]{SkalskiViselter}.

\begin{lemma}[Takai naturality for completely positive maps]
\label{lem:takai}
Let $G$ be a locally compact abelian group, let
$G\curvearrowright_\beta C$, and let $R:C\to C$ be an equivariant completely
positive map. Under the Takai isomorphism
\[
 \Phi_\beta:
 (C\rtimes_\beta G)\rtimes_{\widehat\beta}\widehat G
 \longrightarrow C\otimes\K(L^2G),
\]
one has
\[
 \Phi_\beta\circ\bigl((R\rtimes G)\rtimes\widehat G\bigr)
 =(R\otimes\id_{\K(L^2G)})\circ\Phi_\beta .
\]
\end{lemma}

\begin{proof}
The usual Takai transform is defined first on
$C_c(\widehat G,C_c(G,C))$. Its integral-kernel formula is linear in the
$C$-valued coefficients and uses only scalar Fourier transforms and
$\beta$-translates of those coefficients. The double-crossed extension
applies $R$ to each coefficient. Since $R\beta_s=\beta_sR$, applying $R$
before or after the Takai transform gives the same kernel, with $R$ acting
on its $C$-entry and the identity acting on its compact-operator entry.
The displayed identity follows on the dense double-crossed algebra and hence
everywhere by boundedness. This is the naturality of the standard Takai
duality isomorphism~\cite{ImaiTakai}.
\end{proof}

\begin{lemma}[Stable reflection]\label{lem:stable}
Let $R:C\to C$ be completely positive. If
$R\otimes\id_{\K(H)}$ is nondegenerate for a nonzero Hilbert space $H$, then
$R$ is nondegenerate.
\end{lemma}

\begin{proof}
Let $(e_i)$ be a positive contractive approximate identity of $C$ and let
$(p_j)$ be the net of finite-rank projections on $H$, directed by inclusion
of their ranges. Then $(e_i\otimes p_j)_{(i,j)}$ is an approximate identity
of $C\otimes\K(H)$. Fix $c\in C$ and a rank-one projection $p$ on $H$.
Nondegeneracy of $R\otimes\id$ gives
\[
 \bigl(R(e_i)\otimes p_j\bigr)(c\otimes p)\longrightarrow c\otimes p.
\]
After taking $j$ with $p_jp=p$, this says
$\|R(e_i)c-c\|\to0$. Multiplication on the other side similarly gives
$\|cR(e_i)-c\|\to0$. Thus $R(e_i)\to1$ strictly in $\M(C)$.
\end{proof}

\begin{proposition}[Crossed products reflect nondegeneracy]
\label{prop:reflection}
Let $G$ be a locally compact abelian group and let $R:C\to C$ be an
equivariant completely positive map. If $R\rtimes G$ is nondegenerate, then
$R$ is nondegenerate.
\end{proposition}

\begin{proof}
The map $R\rtimes G$ commutes with the dual action. Since it is
nondegenerate, its crossed-product extension
\[
 (R\rtimes G)\rtimes\widehat G
\]
is nondegenerate by the crossed-product preservation result quoted above.
Lemma~\ref{lem:takai} transports this map to
$R\otimes\id_{\K(L^2G)}$, which is therefore nondegenerate.
Lemma~\ref{lem:stable} now implies that $R$ is nondegenerate.
\end{proof}

\section{\hspace{0.35em}Continuity on a generalized fixed-point algebra}

\begin{lemma}[Continuity of crossed-product semigroups]
\label{lem:crossed-c0}
Let a locally compact group $G$ act on $C$. If $(R_t)_{t\geq0}$ is a
point-norm continuous semigroup of equivariant completely positive
contractions, then $(R_t\rtimes G)_{t\geq0}$ is point-norm continuous on
$C\rtimes G$.
\end{lemma}

\begin{proof}
For $h\in C_c(G,C)$,
\[
 \|(R_t\rtimes G)(h)-h\|_{C\rtimes G}
 \leq \int_G\|R_t(h(s))-h(s)\|\,ds .
\]
The range of $h$ is norm compact. Pointwise convergence $R_t(c)\to c$ and
contractivity imply uniform convergence on norm-compact subsets of $C$.
The right-hand side therefore tends to zero as $t\downarrow0$. Since
$C_c(G,C)$ is dense in the crossed product and the maps are contractions,
the conclusion extends to every element. The semigroup law then gives
continuity at every $t\geq0$.
\end{proof}

\begin{lemma}[Averaging-core transfer]\label{lem:averaging}
Let $(B,\alpha,\phi)$ be a weakly proper $G\rtimes G$-algebra, where $G$
acts on itself by translation and
$\phi:C_0(G)\to\M(B)$ is the structural map. Let $(T_t)_{t\geq0}$ be
nondegenerate completely positive contractions such that:
\begin{enumerate}
  \item $t\mapsto T_t(b)$ is norm continuous for every $b\in B$;
  \item $T_t$ commutes with $\alpha$;
  \item
  \[
   T_t\bigl(\phi(f)b\phi(g)\bigr)
   =\phi(f)T_t(b)\phi(g)
  \]
  for $f,g\in C_c(G)$ and $b\in B$.
\end{enumerate}
Then the restrictions of $T_t$ to the generalized fixed-point algebra
$\Fix(B)$ form a point-norm continuous semigroup.
\end{lemma}

\begin{proof}
Put $B_c=\phi(C_c(G))B\phi(C_c(G))$. For the proper translation action,
the averaging map
\[
 E(a)=\int_G^{\mathrm{st}}\alpha_s(a)\,ds
\]
maps $B_c$ onto a dense subalgebra of $\Fix(B)$. More precisely, for fixed
$f,g\in C_c(G)$ the map
\[
 E_{f,g}:B\longrightarrow\M(B),\qquad
 b\longmapsto E(\phi(f)b\phi(g))
\]
is norm continuous~\cite[Lemma~2.5(4)]{BussEchterhoff}, and the span of its
ranges is dense in the generalized fixed-point algebra
\cite[Proposition~2.2 and Definition~2.8]{BussEchterhoff}.

Nondegeneracy gives a strictly continuous extension of $T_t$ on bounded
subsets of $\M(B)$. Equivariance and the bimodule identity therefore permit
$T_t$ to pass through the strict averaging integral:
\[
 T_t\!\left(E(\phi(f)b\phi(g))\right)
 =E\!\left(\phi(f)T_t(b)\phi(g)\right).
\]
Consequently, norm continuity of $T_t(b)$ and boundedness of $E_{f,g}$ give
point-norm continuity on the dense averaging core. Contractivity extends it
to $\Fix(B)$.
\end{proof}

\section{Proof of the main theorem}

\begin{proof}[Proof of Theorem~\ref{thm:main}]
Set
\[
 B=A\rtimes_\rho\Gamma,
 \qquad \alpha=\widehat\rho^\Psi .
\]
Proposition~4.5 of the source supplies a nondegenerate completely positive
map $\widetilde S:B\to B$ which commutes with $\alpha$, fixes the canonical
copy $\lambda(C_0(\widehat\Gamma))$, and restricts to $S^\Psi$ on
$A^\Psi$~\cite{SkalskiViselter}.

Landstad duality gives a canonical identification
\[
 B=A^\Psi\rtimes_{\rho^\Psi}\Gamma
\]
which preserves the canonical copies of $A^\Psi$ and $\Gamma$. Under this
identification, the multiplicative-domain formula from Proposition~4.5 gives
\[
 \widetilde S=S^\Psi\rtimes\Gamma
\]
on $C_c(\Gamma,A^\Psi)$. Since $\widetilde S$ is nondegenerate,
Proposition~\ref{prop:reflection} shows that $S^\Psi$ is nondegenerate.
This answers Question~4.7.

Now let $(S_t)_{t\geq0}$ be as in Question~4.8 and let
$\widetilde S_t=S_t\rtimes\Gamma$ on $B$. By
Lemma~\ref{lem:crossed-c0}, $(\widetilde S_t)$ is point-norm continuous on
$B$. The triple
\[
 (B,\widehat\rho^\Psi,\lambda:C_0(\widehat\Gamma)\to\M(B))
\]
is a weakly proper
$\widehat\Gamma\rtimes\widehat\Gamma$-algebra, and its generalized
fixed-point algebra is precisely the Landstad algebra $A^\Psi$
\cite[Theorem~4.6]{BussEchterhoff}. Proposition~4.5 of the source gives the
equivariance and bimodule identities required in
Lemma~\ref{lem:averaging}. That lemma therefore shows that the restrictions
$S_t^\Psi$ form a point-norm continuous semigroup on $A^\Psi$. This answers
Question~4.8.
\end{proof}

\section{Verification and novelty status}

The proof has three independent consistency checks.
\begin{itemize}
  \item If $\Gamma$ is discrete, $\widehat\Gamma$ is compact and the
  averaging-core proof reduces to the compact fixed-point argument in the
  earlier partial packet.
  \item If the cocycle is trivial, the Landstad algebra is the undeformed
  algebra and both conclusions reduce to their hypotheses.
  \item No conclusion is drawn from convergence in the ambient strict
  topology. Question~4.7 is transferred to a stabilization, while
  Question~4.8 is proved in norm on a dense generalized fixed-point core.
\end{itemize}

A bounded novelty search on 23 July 2026 used the exact wording of
Questions~4.7 and 4.8, the source title and authors, citation searches for
arXiv:2311.04630, and combinations of ``Takai duality'', ``crossed-product
completely positive map'', ``generalized fixed-point algebra'',
``nondegenerate'', and ``point-norm continuous''. No paper explicitly
answering either question, and no prior statement of the two arguments above,
was found. Buss and Echterhoff's paper predates the questions and supplies
the averaging-core machinery, but does not state these consequences. This is
a bounded search, not an exhaustive novelty certification.

\medskip
\noindent\textbf{Human-review recommendation.}
Review Proposition~\ref{prop:reflection}, especially Takai naturality for
completely positive maps, and the identification of the Landstad algebra with
the generalized fixed-point completion used in Lemma~\ref{lem:averaging}.
These are the only two non-elementary bridges in the proof.

\begin{thebibliography}{9}
\bibitem{SkalskiViselter}
A.~Skalski and A.~Viselter,
\emph{Convolution semigroups on Rieffel deformations of locally compact
quantum groups}, arXiv:2311.04630; Lett. Math. Phys. \textbf{114} (2024), 52.

\bibitem{BussEchterhoff}
A.~Buss and S.~Echterhoff,
\emph{Universal and exotic generalized fixed-point algebras for weakly
proper actions and duality}, arXiv:1304.5697.

\bibitem{ImaiTakai}
S.~Imai and H.~Takai,
\emph{On a duality for $C^*$-crossed products by a locally compact group},
J. Math. Soc. Japan \textbf{30} (1978), 495--504.
\end{thebibliography}

\end{document}
